\section{Multimodal Foundation Agents：VIMA、RT-2 与 RoboCat}

视频路线扩大数据，另一条路线则统一接口。\term{multimodal prompting（多模态提示）}指 prompt 中交错出现文本、图像、目标帧、视频示范或其他感知 token，模型据此输出 action；它不是给语言 prompt 附一张装饰图片，而是让视觉对象、约束和示范直接成为任务规范。

\subsection{统一 Text、Image、Video、Audio 与 Action}

一个 multimodal foundation model 试图把不同输入映射到共享序列空间，再依据上下文生成语言或 action。对 agent 而言，优势是任务可由“把这个物体放到像这张图的位置”或“一次示范”表达；挑战是每种 modality 的 tokenization、时序、精度和 safety 需求不同。

\lecturefigure{slide-50-multimodal-foundation-model.jpg}{统一文本、图像、视频、声音与 action 的 foundation model 愿景}{官方录像恢复幻灯片 V050；00:43:03}

\begin{knowledgebox}{读图：统一 token space 不等于统一物理意义}
图中 modality 都进入一个模型，但文本 token、视觉 patch 和 robot action 的误差代价不同。先看输入怎样编码，再看输出是否受到 action-space constraint；统一 architecture 有利于 knowledge transfer，却仍需要 modality-specific encoder、controller 与 safety validation。
\end{knowledgebox}

\lecturefigure{slide-51-multimodal-prompting-robots.jpg}{机器人任务可以由交错的视觉与语言 prompt 定义}{官方录像恢复幻灯片 V051；00:43:48}

Prompt 可包含目标物体图片、位置示意、禁止区域和自然语言关系。读图时先区分“指代哪个对象”“目标状态是什么”“动作必须满足哪些 constraint”，再观察 robot trajectory。多模态 prompt 提升 specification bandwidth，但也引入 ambiguity；模型必须 grounding 到当前 scene，而不是匹配训练中的外观模板。

\subsection{VIMA：Visual Manipulation with Multimodal Prompts}

VIMA 把视觉对象 token 与文字 token 交错输入 Transformer，并 autoregressively 预测 robot action。它的重点不是单个 manipulation benchmark，而是用同一 model/interface 表达 visual goal reaching、novel concept grounding、one-shot demonstration 和 constraint satisfaction 等不同任务族。

\lecturefigure{slide-52-vima-architecture.jpg}{VIMA 将 multimodal prompt、observation 与 action 统一建模}{官方录像恢复幻灯片 V052；00:44:33}

\begin{knowledgebox}{读图：prompt encoder 与 action decoder 的连接}
先看 prompt 中图像和文本如何组成 token sequence，再看当前 observation 如何加入上下文，最后看 action decoder 输出 pose 或 control token。VIMA 的共享 Transformer 促进任务间迁移，但 benchmark 中的 tabletop world、object set 和 action representation 仍限定其证据范围。
\end{knowledgebox}

若 prompt token 为 $p_{1:m}$、历史 observation 为 $o_{1:t}$、action sequence 为 $a_{1:t}$，autoregressive objective 可写成：
\[
\mathcal{L}_{\mathrm{VIMA}}(\theta)
=-\sum_{t=1}^{T}\log p_\theta\!\left(a_t\mid p_{1:m},o_{1:t},a_{<t}\right).
\]
其中 $p_{1:m}$ 可以交错包含文本与视觉对象，$o_{1:t}$ 是当前及历史 observation，$a_{<t}$ 是先前动作，$T$ 是 episode horizon，$\theta$ 是模型参数。公式表达统一 conditional sequence modeling，并不意味着所有 action error 都与语言 token error 等价。

\subsection{四类 Prompt 能力的证据边界}

下面四页分别测试视觉目标、临时概念、一次示范和约束。应把它们视为 task specification 的不同形式，而不是四个互不相关的 demo；读图时持续检查 prompt 提供了什么信息、模型需要 grounding 哪个对象、评价是否只在受控 tabletop distribution 内进行。

\lecturefigure{slide-53-vima-visual-goal-reaching.jpg}{VIMA：由目标图像指定希望达到的视觉状态}{官方录像恢复幻灯片 V053；00:44:51}

视觉目标减少了语言描述空间关系的负担。先比较初始 scene 与目标 image，再检查动作后对象相对位置；成功说明模型能把目标外观映射为 manipulation sequence，但无法单独证明它理解了更广泛的物理因果或真实世界遮挡。

\lecturefigure{slide-54-vima-concept-grounding.jpg}{VIMA：在 prompt 内 grounding 临时定义的新概念}{官方录像恢复幻灯片 V054；00:45:03}

“wug”“blicket”等无既定语义的词通过 prompt 中的视觉对应被临时定义。读图时先识别 token 与对象的 binding，再看模型是否把关系迁移到 action；这支持 in-context concept grounding，而不是证明模型预训练时已经见过这些词的含义。

\lecturefigure{slide-55-vima-one-shot-demonstration.jpg}{VIMA：从一次视觉示范推断任务并复现}{官方录像恢复幻灯片 V055；00:45:15}

一次示范提供的是 trajectory-level specification。先看 demonstration 中对象变化，再看 test scene 是否具有新位置或新实例；成功要求模型抽取关系而非复制绝对坐标。单次 demo 能否推广仍依赖任务族、摄像机和 object distribution 的相似性。

\lecturefigure{slide-56-vima-constraint-satisfaction.jpg}{VIMA：在视觉约束下完成 manipulation}{官方录像恢复幻灯片 V056；00:45:48}

Constraint prompt 告诉 agent 哪个区域、对象或路径不能触碰。读图时不只检查最终目标，还要检查 trajectory 是否违反中间约束；这说明 agent evaluation 必须包含过程安全，而不是只看 final frame。受控桌面约束仍不同于真实机器人的 force、collision 与 human safety。

\teachervoice{讲者说 VIMA “tokenizes everything”，用 interleaved image--text tokens 统一过去需要专门 pipeline 的任务。这句话应理解为 modeling interface 的统一，而不是物理控制问题消失；视觉 grounding、action decoder 与环境约束仍各自承担责任。}

\subsection{RT-2 与 RoboCat：扩大 Robot Data 与 Transfer}

课程最后用 RT-2 与 RoboCat 展示同一 recipe 的后续实例。RT-2 把 web-scale vision-language knowledge 与 robot demonstration co-fine-tune，并把 action 表示为可生成 token；RoboCat 训练统一 policy 跨任务和 robot embodiment。两者扩大了 transfer 证据，但仍依赖具体数据混合与 action support。

\lecturefigure{slide-57-rt2-robocat.jpg}{RT-2 与 RoboCat 延伸 multimodal generalist robot 路线}{官方录像恢复幻灯片 V057；00:46:12}

\begin{knowledgebox}{术语消化：三种 generalist-agent 接口}
\begin{tabularx}{\textwidth}{@{}L{0.17\textwidth}L{0.25\textwidth}X@{}}
\toprule
系统 & 主要接口 & 核心机制与边界 \\
\midrule
Voyager & JavaScript code-as-action & 长 horizon、可组合、依赖 symbolic API \\
Eureka & LLM-generated reward code & 用 RL 学连续控制、依赖 simulator 与 reward audit \\
VIMA / RT-2 / RoboCat & Multimodal token 到 robot action & 共享表示与数据迁移、依赖 action grounding 和 embodiment coverage \\
\bottomrule
\end{tabularx}
\end{knowledgebox}

\teachervoice{讲者把 RT-2 解释为先在 internet-scale data 上预训练，再用 human-collected robot demonstrations fine-tune；RoboCat 则展示单一 policy 跨机器人与任务。课堂把它们当作 promising follow-up，而没有宣称 generalist robotics 已经完成。}

\subsection{本章小结}

Multimodal prompting 扩大了任务 specification：文字、目标图像、临时概念、示范和 constraint 都可进入统一上下文。VIMA 证明多类 tabletop task 可由同一 Transformer/action interface 表达，RT-2 与 RoboCat 进一步扩大 web prior、robot data 和 embodiment transfer。统一模型带来共享表示，但 grounding、control、safety 与 evidence boundary 仍必须单独检查。

